\chapter{Fazit}

Ziel dieser Bachelorarbeit war es, eine VirtIO-Treiberbasis in D3OS zu integrieren, indem die Rust-Crate virtio-drivers so angebunden wird, dass VirtIO-Geräte künftig über ein einheitliches und wiederverwendbares Treibermodell nutzbar sind. Dieses Ziel wurde erreicht: Die betriebssystemspezifischen Aufgaben wie Speicherbereitstellung für DMA, Adressübersetzung und die Anbindung an PCI wurden so gekapselt, dass die generische VirtIO-Logik sauber darauf aufsetzen kann. Dadurch entsteht eine tragfähige Grundlage, um weitere VirtIO-Geräte mit deutlich geringerem Implementierungsaufwand zu integrieren, statt pro Gerät eine isolierte Treiberlösung zu entwickeln.\\
\\
Auf dieser Basis wurden zwei konkrete Funktionsbereiche umgesetzt und evaluiert. Für die VirtIO-GPU wurde nicht nur die grundlegende Grafikfunktionalität in D3OS verfügbar gemacht, sondern der Treiber außerdem um Mechanismen erweitert, die das Senden von VirGL/3D-Command-Streams ermöglichen. Der durchgeführte VirGL-Smoke-Test liefert dabei den End-to-End-Nachweis, dass der Befehlsfluss vom Gast bis zur Host-Verarbeitung funktioniert und Ergebnisse korrekt dargestellt werden, während gleichzeitig auch dynamische Größenänderungen (Resize) implementiert und in Demo-Programmen erprobt wurden. Die quantitative Evaluation der Grafikleistung offenbarte jedoch eine starke Abhängigkeit von der verwendeten Virtualisierungsplattform: Während ein nativer Linux-KVM-Host eine vernachlässigbare Latenz und hohe Bildraten ermöglichte, führte die Ausführung unter WSL2 (Hyper-V) bei steigenden Auflösungen zu signifikanten Leistungseinbußen Für VirtIO-Sound wurde eine funktionierende PCM-Playback-Pipeline realisiert und über eine Beispielwiedergabe validiert, sodass Audioausgabe in der VM praktisch nutzbar ist.\\
\\
Trotz der erfolgreichen Integration unterliegt die Arbeit gewissen Limitationen, insbesondere im Bereich der Leistungsmessung. Die notwendige Deaktivierung von Interrupts zur Synchronisation in kritischen Treiberabschnitten beeinträchtigte den System-Timer, was eine nachträgliche Korrektur der Zeitbasis bei Performance-Tests erforderte. Für zukünftige Weiterentwicklungen empfiehlt sich daher die Umstellung der Zeitmessung auf den CPU-internen Time Stamp Counter (TSC) sowie eine verfeinerte, selektive Interrupt-Maskierung, oder die Nutzung von Polling. Zusammenfassend lässt sich jedoch festhalten, dass die Integration von VirtIO-Treibern die Portabilität und Testbarkeit von D3OS entscheidend verbessert und eine solide Basis für zukünftige Projekte darstellt.